\section*{Chapitre \ref{ch:pl:distributed-compute-and-schedulers} --- Calcul distribué et ordonnanceurs}

\begin{solution}{exo:pl:distributed-compute-and-schedulers:1}
La règle du plus long d'abord place 3 et 3 sur les deux cœurs, puis 2 et 2, puis le dernier 2 sur le premier cœur libre~: 7. Le meilleur est 6 (3 et 3 sur un même cœur, 2,
2 et 2 sur l'autre). Le rapport $7/6$ égale la borne de Graham $4/3 - 1/(3m)$ pour $m = 2$~: l'exemple est le pire cas.
\end{solution}

\begin{solution}{exo:pl:distributed-compute-and-schedulers:2}
Trente fois quinze minutes, soit 7.5~heures, sur un nœud normal~; quatre fois cette durée, soit 30~heures, sur le nœud lent.
\end{solution}

\begin{solution}{exo:pl:distributed-compute-and-schedulers:3}
Le travail total est de 3\,649 heures-cœur~; réparti sur 384 cœurs, cela fait 9.50~heures, plus que la tâche la plus longue (7.5~heures sur un nœud normal) et
que toute chaîne (le balayage à plat n'en a aucune). La borne retient la plus grande des trois.
\end{solution}

\begin{solution}{exo:pl:distributed-compute-and-schedulers:4}
Une tâche saine s'exécute pendant sa durée réelle, qui ne dépasse 2.5~fois son estimation que lorsque l'erreur de l'estimation $e^{0.3z}$ est inférieure à $1/2.5$, soit
$z > 3.05$~: une probabilité de 0.11~\%. Avec un facteur de 1.5, le seuil est franchi dès que $z > 1.35$, soit environ 9~\% des tâches~: des centaines de copies de secours inutiles,
chacune prenant un cœur à la file d'attente et gaspillant du travail.
\end{solution}

\begin{solution}{exo:pl:distributed-compute-and-schedulers:5}
Le \omterm{def:pl:distributed-compute-and-schedulers:stealing}{vol de tâches} occupe chaque cœur tant qu'il reste du travail, mais il prend les tâches dans l'ordre où elles ont été distribuées~; les tâches longues, soumises en dernier, démarrent
tard et finissent tard, comme sous le premier arrivé, premier servi. L'équilibrage n'est pas l'ordre~: seule une politique qui lance d'abord les tâches les plus longues évite la queue.
\end{solution}

\begin{solution}{exo:pl:distributed-compute-and-schedulers:6}
Dans cette exécution, il se termine plus tôt, à 16.68~heures au lieu de 19.79~: les 500 tâches de la deuxième équipe ne représentent qu'environ 5~\% du travail, et la première
équipe est fixée par l'endroit et le moment où ses quarante tâches longues s'exécutent, ce que l'entrelacement modifie --- ici à son avantage, mais
peut-être pas dans une autre exécution.
\end{solution}

\begin{solution}{exo:pl:distributed-compute-and-schedulers:7}
Le premier arrivé, premier servi s'améliore, de 19.79 à 14.41~heures (13.69 avec copies de secours)~: les tâches longues démarrent aussitôt. Il reste derrière le plus long d'abord,
dont le résultat reste inchangé à 11.80 (10.02 avec copies de secours), car il trie selon l'estimation et non selon la date de soumission~: les autres tâches à queue lourde sont toujours dans
un ordre aléatoire sous le premier arrivé, premier servi.
\end{solution}

\begin{solution}{exo:pl:distributed-compute-and-schedulers:8}
Le taux d'utilisation sur une nuit ne dit rien de la manière dont la capacité a été utilisée~: le balayage a immobilisé 475.0 heures-nœud pour 3\,649 heures-cœur de travail à cause de
l'ordre de ses tâches, et il est resté presque inactif pendant trois heures et demie. Sur la moitié de la grappe, la borne inférieure double pour atteindre 19~heures~; le remède est
l'ordonnancement (le plus long d'abord, copies de secours), qui finit près de la borne sur la grappe telle qu'elle est.
\end{solution}

\begin{solution}{pb:pl:distributed-compute-and-schedulers:1}
\begin{enumerate}
\item Chaque backtest lit ses entrées et écrit son résultat sans communiquer avec les autres.
\item Sur 24 petits backtests du chapitre~11 exécutés sur cet ordinateur portable via l'exécuteur local~: la dispersion des logarithmes des durées, 0.76.
\item 3\,649 heures-cœur~; 9.50~heures sur 384 cœurs.
\item À la fin, parce qu'une grille de paramètres soumet ses valeurs les plus grandes en dernier.
\item Longueur intrinsèque, nœud lent et échec relancé depuis le début.
\item Au plus $2 - 1/m$ fois la \omterm{def:pl:distributed-compute-and-schedulers:makespan}{durée totale} optimale pour tout ordre~; au plus $4/3 - 1/(3m)$ pour le plus long d'abord.
\item Estimations de durée.
\item Il équilibre la charge mais conserve l'ordre de distribution, si bien que les tâches longues démarrent encore en dernier.
\item Il sert d'abord l'équipe qui utilise moins que sa part, de sorte qu'un petit balayage tardif ne soit pas mis en file derrière un grand.
\item Les copies de secours aident quand une tâche est lente à cause de l'endroit où elle tourne (un nœud lent)~; elles gaspillent quand les estimations sont mauvaises et que le seuil se déclenche sur des tâches saines.
\item Premier arrivé, premier servi 19.79 (20.32 avec copies de secours)~; le plus long d'abord 11.80 (10.02)~; \omterm{def:pl:distributed-compute-and-schedulers:stealing}{vol de tâches} 16.33 (16.40)~; le plus long d'abord sans le
nœud lent 10.32.
\item Environ trois heures et demie~: moins d'un dixième des cœurs étaient occupés de 16.3~heures à 19.8.
\item 5.5~\% (10.02 contre 9.50), à cause du nœud lent et de l'erreur d'estimation~; sans le nœud lent mais sans copies de secours, 10.32.
\item La moitié de ses tâches sont terminées 0.59~heure après la soumission, au lieu de 8.28~heures.
\item Le graphe calcule les variables une seule fois (3\,514 au lieu de 6\,922 heures-cœur, 11.4 au lieu de 19.3~heures)~; le cache évite les cent
étapes de calcul des variables du balayage suivant (3\,475 heures-cœur, 10.3~heures).
\item \textbf{Résultat nommé.} Sur 384 cœurs, le balayage de 10\,000 tâches dure 19.79~heures et 475.0 heures-nœud en premier arrivé, premier servi, 11.80~heures et 283.3
heures-nœud pour le plus long d'abord, 16.33~heures et 391.9 heures-nœud avec le \omterm{def:pl:distributed-compute-and-schedulers:stealing}{vol de tâches}~; les copies de secours dès 2.5~fois l'estimation font gagner 1.8~heure au plus long d'abord et coûtent 0.5~heure au premier arrivé,
premier servi~; le meilleur ordonnancement, le plus long d'abord avec copies de secours, finit 5.5~\% au-dessus de la borne inférieure de 9.50~heures.
\item Le plus long d'abord selon les estimations par défaut, copies de secours avec un seuil fixé d'après l'erreur des estimations, \omterm{def:pl:distributed-compute-and-schedulers:fairshare}{ordonnancement à parts équitables} entre les équipes.
\item Une estimation de la durée de chaque tâche (ou des paramètres qui la prédisent), et des clés de cache pour les étapes partagées.
\item À partir de la borne inférieure du balayage type et de l'échéance~: cœurs = travail / (échéance $-$ marge de la tâche la plus longue), puis vérifier avec le
simulateur sous la politique réellement utilisée.
\item Ses dernières tâches~: leur durée, l'endroit où elles s'exécutent et le moment où elles démarrent.
\end{enumerate}
\end{solution}

\begin{solution}{iq:pl:distributed-compute-and-schedulers:1}
Il faut trouver pourquoi elles sont longues~: une longueur intrinsèque (les lancer en premier, avec le plus long d'abord selon les estimations), une machine lente (copies de secours sur un autre nœud),
ou des échecs (points de contrôle, reprises)~; il faut ensuite comparer la \omterm{def:pl:distributed-compute-and-schedulers:makespan}{durée totale} à la borne inférieure.
\iqlookfor{Diagnostiquer la queue avant d'ajouter de la capacité.}
\end{solution}

\begin{solution}{iq:pl:distributed-compute-and-schedulers:2}
Lancer une seconde copie d'une tâche en retard, sur une autre machine, en retenant la première qui se termine. C'est une mauvaise idée quand les tâches sont en retard parce que
elles sont longues plutôt qu'à cause de l'endroit où elles tournent, quand les estimations sont trop mauvaises pour trancher, ou quand le travail a des effets de bord que deux copies
dupliqueraient.
\iqlookfor{Seuils fixés d'après l'erreur d'estimation~; tâches idempotentes.}
\end{solution}

\begin{solution}{iq:pl:distributed-compute-and-schedulers:3}
\omterm{def:pl:distributed-compute-and-schedulers:fairshare}{Ordonnancement à parts équitables} avec des cibles par équipe, de sorte qu'une équipe en deçà de sa part soit servie en premier~; des limites de cœurs simultanés par utilisateur~; et une priorité pour
les travaux interactifs courts sur les longs balayages.
\iqlookfor{Des parts, et non le premier arrivé, premier servi.}
\end{solution}

\begin{solution}{iq:pl:distributed-compute-and-schedulers:4}
Sous forme de \omterm{def:pl:distributed-compute-and-schedulers:graph}{graphe de tâches}~: les calculs de variables comme étapes partagées, avec des clés de cache adressées par le contenu, et les backtests qui en dépendent~; puis ordonnancer
le graphe, les tâches estimées les plus longues d'abord.
\iqlookfor{Graphe, clés de cache, réutilisation d'un balayage à l'autre.}
\end{solution}

\begin{solution}{iq:pl:distributed-compute-and-schedulers:5}
\omterm{def:pl:distributed-compute-and-schedulers:graph}{Graphes de tâches} avec demandes de ressources et estimations~; une file avec \omterm{def:pl:distributed-compute-and-schedulers:fairshare}{ordonnancement à parts équitables} entre équipes et ordre du plus long d'abord au sein d'un balayage~; reprises
avec points de contrôle~; copies de secours avec un seuil calibré~; un cache adressé par le contenu~; comptabilité des heures-nœud par équipe~; et un simulateur de la
grappe pour tester les politiques avant leur déploiement.
\iqlookfor{Politique choisie par la mesure et la simulation, et non par défaut.}
\end{solution}
